\subsection{Excluding \texorpdfstring{$C_4\times C_{20}$}{C4 x C20}}
\begin{lemma}\label{lem:decic420-gonality}
The gonality of $X_1(4,20)$ over $\mathbb F_9$ is at least $7$.
\end{lemma}
\begin{proof}
The numbers of closed places in degrees one, two and three are
$24$, $108$ and $320$. These counts agree with the Frobenius
polynomial of $J_H(80)$, $H=\langle21\rangle$, after quadratic
constant-field extension.
A map to $\mathbb P^1$ of degree at most $6$ would have a rational
fiber containing three distinct rational points: there are $24$
rational points on the curve and only ten rational fibers.
After adding rational points if necessary, we obtain an effective
divisor $A+T$ of degree $6$, where $\deg A=3$ and $T$ consists
of three distinct rational points, with $\ell(A+T)\ge2$.

There are
\[
 \binom{26}{3}+108\cdot24+320=5512
\]
effective degree-three divisors $A$, and $\binom{24}{3}=2024$
choices of $T$. We compute the canonical vanishing conditions
for every $A$. Each space $H^0(K-A)$ has dimension $22$.
The three additional jet conditions at the points of $T$ have
rank three, including when $A$ and $T$ intersect. Hence
$\ell(A+T)=1$ in all
\[
 5512\cdot2024=11156288
\]
cases, a contradiction.
The exported local differential matrices include every required
multiplicity. Independent arithmetic over $\mathbb F_9$
reconstructs all $5512$ kernels and their projected row spaces,
and repeats every three-column rank test.
\end{proof}

\begin{lemma}\label{lem:decic-cusp-correction}
Let $X/F$ have rank-zero Jacobian and $F$-gonality greater than
$d+1$, where $d$ is odd. Let $v$ be an unramified odd prime
of good reduction. Suppose that an effective noncuspidal
$F$-rational divisor $D$ of degree $d$ has reduction $B+\bar c$,
where $c$ is an $F$-rational cusp. If an $F$-defined involution
$g$ preserves the cusp boundary, fixes $B$, and acts freely on
the noncuspidal locus, then no such $D$ exists.
\end{lemma}
\begin{proof}
The torsion class $[D-gD-c+gc]$ reduces to zero. Odd unramified
torsion injectivity gives $D+gc\sim gD+c$.
Both effective divisors have degree $d+1$, so the gonality
hypothesis gives equality. Separating cuspidal and noncuspidal
support yields $D=gD$. Freeness of the involution forces
$\deg D$ to be even, a contradiction.
\end{proof}

\begin{proposition}\label{prop:decic420}
No elliptic curve over a decic number field contains $C_4\times C_{20}$.
\end{proposition}
\begin{proof}
Put $F=\Q(i)$, $X=X_1(4,20)$ and $J=\operatorname{Jac}(X)$.
A proposed subgroup gives an effective noncuspidal $F$-rational
divisor $D$ of degree $5$. The curve has genus $25$, and
$J(F)$ has rank zero by \cite[Theorem~4.1]{DS}.
Lemma~\ref{lem:decic420-gonality} and good reduction give
$\operatorname{gon}_F X\ge7$.

Reduce at $3$, whose residue field is $\mathbb F_9$.
The cusp signature is $1^{24},2^{12}$, and the noncuspidal
counts in degrees one through five are
\[
 (0,96,320,888,11904).
\]
Lemma~\ref{lem:decic-cuspidal-reduction} excludes entirely
cuspidal reduction. All rational residue cusps are globally
$F$-rational, as verified by the complete primitive Tate-cusp
action. The other global cusp orbits have degree four.

For a split prime $q$ write $A_q=T_q-1-q\langle q\rangle$.
Eichler--Shimura and odd torsion injectivity give $A_qJ(F)=0$.
The boundary formula shows that $A_{29}$ kills every closed
cuspidal residue divisor, since $29\equiv9\pmod{20}$.
The operators $A_{13}$ and $A_{17}$ kill every rational cusp.
Thus the following principality tests cover all possible
noncuspidal support patterns of the reduction.
\[
\begin{array}{c|r|l|r}
\text{support degrees}&\text{representatives tested}&
\text{Hecke primes}&\text{remaining supports}\\\hline
{[2]}&2&29&0\\
{[3]}&4&29&0\\
{[4]}&12&29,13,17&2\\
{[5]}&124&13&0\\
{[2,2]}&192&29,13,17&2\\
{[2,3]}&640&13&0
\end{array}
\]
Here the full pairing-preserving level group is
\[
 G=(\operatorname{SL}_2(\mathbb Z/4\mathbb Z)
       \times(\mathbb Z/5\mathbb Z)^*)/\langle(-I,-1)\rangle,
 \qquad |G|=96.
\]
For a pair we normalize its quadratic support and retain every
second support, including repeated quadratic places.
All point counts and every orbit partition are checked
independently. The twelve quintic jobs cover all $124$ orbits;
the eight mixed-support jobs cover all $832$ pairs.

Each remaining support has degree four, so its complement is
one rational cusp. Write the level structure as $(P_4,Q_{20})$
and put $Q_4=5Q_{20}$, $P_5=4Q_{20}$.
For the two quartic supports the stabilizers have order eight
and contain the involution with $4$-level matrix $-I$ and
trivial $5$-component. For the two quadratic pairs the stabilizers
have order four and contain, respectively, the involutions
\[
 \begin{pmatrix}3&2\\2&3\end{pmatrix},\qquad
 \begin{pmatrix}1&2\\0&1\end{pmatrix}
\]
on $(P_4,Q_4)$, again with trivial $5$-component.
The exact action of every element of $G$ verifies these stabilizers.

All three involutions act freely on the noncuspidal locus.
Indeed, they fix $P_5$, and an automorphism of an elliptic curve
fixing a point of order five is the identity. Their nontrivial
action on the full $4$-level structure therefore admits no fixed
point. This also covers CM points.
Lemma~\ref{lem:decic-cusp-correction}, with $d=5$, excludes all
four remaining supports and proves the proposition.
\end{proof}
